\chapter{Documentation des API REST}\label{chap:rest}

\section{Conventions}

Les trois services REST partagent les conventions suivantes.

\begin{itemize}
  \item \textbf{Ressources et verbes.} Les URI désignent des ressources au pluriel (\code{/candidatures}, \code{/stages/\{id\}/evaluations}), versionnées sous \code{/api/v1}. \code{GET} consulte, \code{POST} crée ou déclenche une action métier, \code{PUT} remplace, \code{PATCH} modifie partiellement (le statut), \code{DELETE} supprime ou annule.
  \item \textbf{Format.} JSON en UTF-8 dans les deux sens ; dates au format ISO~8601.
  \item \textbf{En-têtes.} \code{X-API-Key} (obligatoire, une clé par client), \code{X-Correlation-Id} (facultatif, renvoyé dans la réponse), \code{X-Acteur} (auteur métier d'une action, par exemple \code{entreprise:1}, inscrit dans les historiques).
  \item \textbf{Contrats.} Chaque service sert son contrat OpenAPI~3 (\code{/openapi.yaml}) et une interface Swagger UI (\code{/docs/}) fonctionnant sans Internet. Un test vérifie que chaque route du code figure dans le contrat et réciproquement.
\end{itemize}

\begin{table}[H]
\centering\small
\begin{tabularx}{\linewidth}{@{}l l>{\raggedright\arraybackslash}X@{}}
\toprule
\textbf{Code} & \textbf{Code d'erreur} & \textbf{Usage} \\ \midrule
200 & – & Consultation, modification, décision réussies. \\
201 & – & Ressource créée (candidature, notification, stage, compte rendu, évaluation). \\
202 & – & Soumission traitée par le processus : enregistrée ou rejetée (le résultat est dans le corps). \\
204 & – & Annulation ou suppression réussie, sans corps. \\
400 & \code{REQUETE\_INVALIDE} & Requête mal formée : JSON illisible, paramètre de type incorrect. \\
401 & \code{NON\_AUTHENTIFIE} & Clé d'API absente ou invalide. \\
403 & \code{ACCES\_INTERDIT} & Clé valide, mais client non autorisé pour cette opération. \\
404 & \code{*\_INTROUVABLE} & Ressource inexistante (\code{CANDIDATURE\_INTROUVABLE}, \code{STAGE\_INTROUVABLE}…). \\
409 & selon la règle & Conflit avec une règle de gestion (R5, R6, R8, R9, R10). \\
422 & \code{DONNEES\_INVALIDES} & Données bien formées mais invalides (champ manquant, note hors de 0 à 20…). \\
503 & \code{SERVICE\_*\_INDISPONIBLE} & Un service dépendant (\code{SERVICE\_OFFRES\_INDISPONIBLE}…) ou la base (\code{BASE\_INDISPONIBLE}) ne répond pas. \\
\bottomrule
\end{tabularx}
\caption{Codes HTTP utilisés}\label{tab:codes-http}
\end{table}

Toute erreur a le même corps JSON : \code{code} (stable), \code{message} (en français), \code{details} (liste, par exemple un message par champ invalide), \code{timestamp} et \code{path}. Le listing~\ref{lst:erreur-409} en donne un exemple réel.

\section{Points d'accès}

Comme pour SOAP, les tableaux sont produits à partir des contrats OpenAPI (script \code{generer\_tables.py}). Un point d'interrogation marque un paramètre facultatif ; la colonne « Corps » nomme le schéma OpenAPI attendu.

\subsection{S4 Gestion des candidatures}
S4 distingue la \emph{soumission} (la demande qui entre dans le processus) de la \emph{candidature} (créée seulement si la passerelle de décision l'accepte). Les points d'accès marqués « processus » sont appelés par Bonita : contrôles R5 et R6, enregistrement, rejet, journal des étapes.
{\small% Généré par rapport/outils/generer_tables.py à partir de rest-candidatures/src/main/resources/openapi.yaml. Ne pas modifier.
\begin{longtable}{@{}>{\raggedright\arraybackslash}p{0.06\linewidth}>{\raggedright\arraybackslash}p{0.40\linewidth}>{\raggedright\arraybackslash}p{0.16\linewidth}>{\raggedright\arraybackslash}p{0.14\linewidth}>{\raggedright\arraybackslash}p{0.12\linewidth}@{}}
\caption{Points d'accès REST de S4 Candidatures (base \texttt{/api/v1}, 23 opérations)}\label{tab:rest-candidatures}\\
\toprule \textbf{Verbe} & \textbf{URI et rôle} & \textbf{Paramètres} & \textbf{Corps} & \textbf{Codes} \\ \midrule
\endfirsthead
\toprule \textbf{Verbe} & \textbf{URI et rôle} & \textbf{Paramètres} & \textbf{Corps} & \textbf{Codes} \\ \midrule
\endhead
\bottomrule
\endlastfoot
\textsc{post} & \texttt{/\allowbreak{}soumissions}\newline{\footnotesize Déposer une candidature (déclenche le processus)} & – & \texttt{NouvelleSoumission} & 202, 400, 401, 403, 422, 503 \\ \midrule
\textsc{get} & \texttt{/\allowbreak{}soumissions/\allowbreak{}\{id\}}\newline{\footnotesize Consulter une soumission et ses étapes} & \texttt{id} & – & 200, 401, 404 \\ \midrule
\textsc{get} & \texttt{/\allowbreak{}soumissions/\allowbreak{}\{id\}/\allowbreak{}etapes}\newline{\footnotesize Étapes du processus déjà exécutées (vue BPMN en direct)} & \texttt{id} & – & 200, 401, 404 \\ \midrule
\textsc{post} & \texttt{/\allowbreak{}soumissions/\allowbreak{}\{id\}/\allowbreak{}etapes}\newline{\footnotesize Signaler une étape exécutée par le processus Bonita} & \texttt{id} & \texttt{NouvelleEtape} & 201, 400, 401, 403, 404, 422 \\ \midrule
\textsc{post} & \texttt{/\allowbreak{}soumissions/\allowbreak{}\{id\}/\allowbreak{}rejet}\newline{\footnotesize Enregistrer le rejet décidé par le processus Bonita (R1 à R7)} & \texttt{id} & \texttt{Rejet} & 200, 400, 401, 403, 404, 409, 422 \\ \midrule
\textsc{post} & \texttt{/\allowbreak{}soumissions/\allowbreak{}\{id\}/\allowbreak{}erreur}\newline{\footnotesize Signaler l'interruption du processus Bonita (service indisponible)} & \texttt{id} & \texttt{JSON} & 200, 400, 401, 403, 404, 422 \\ \midrule
\textsc{get} & \texttt{/\allowbreak{}processus/\allowbreak{}mode}\newline{\footnotesize Mode d'orchestration en service (bonita ou direct)} & – & – & 200, 401 \\ \midrule
\textsc{get} & \texttt{/\allowbreak{}etudiants/\allowbreak{}\{id\}/\allowbreak{}soumissions}\newline{\footnotesize Soumissions d'un étudiant, y compris rejetées} & \texttt{id} & – & 200, 401, 422 \\ \midrule
\textsc{post} & \texttt{/\allowbreak{}candidatures}\newline{\footnotesize Enregistrer une candidature (appelé par le processus après la passerelle de décision)} & – & \texttt{NouvelleCandidature} & 201, 400, 401, 403, 404, 409, 422, 503 \\ \midrule
\textsc{get} & \texttt{/\allowbreak{}candidatures}\newline{\footnotesize Lister les candidatures (filtres combinables)} & \texttt{statut?}, \texttt{etudiantId?}, \texttt{offreId?}, \texttt{entrepriseId?} & – & 200, 400, 401 \\ \midrule
\textsc{get} & \texttt{/\allowbreak{}candidatures/\allowbreak{}verification-doublon}\newline{\footnotesize Règle R5 - candidature non annulée existante pour ce couple étudiant et offre} & \texttt{etudiantId}, \texttt{offreId} & – & 200, 400, 401 \\ \midrule
\textsc{get} & \texttt{/\allowbreak{}candidatures/\allowbreak{}\{id\}}\newline{\footnotesize Consulter une candidature} & \texttt{id} & – & 200, 401, 404 \\ \midrule
\textsc{delete} & \texttt{/\allowbreak{}candidatures/\allowbreak{}\{id\}}\newline{\footnotesize Annuler une candidature (règle R8)} & \texttt{id} & – & 204, 401, 404, 409 \\ \midrule
\textsc{get} & \texttt{/\allowbreak{}candidatures/\allowbreak{}\{id\}/\allowbreak{}historique}\newline{\footnotesize Historique des changements de statut (date, acteur, commentaire)} & \texttt{id} & – & 200, 401, 404 \\ \midrule
\textsc{get} & \texttt{/\allowbreak{}candidatures/\allowbreak{}\{id\}/\allowbreak{}etapes}\newline{\footnotesize Étapes du processus de la candidature (vue BPMN en direct)} & \texttt{id} & – & 200, 401, 404 \\ \midrule
\textsc{patch} & \texttt{/\allowbreak{}candidatures/\allowbreak{}\{id\}/\allowbreak{}statut}\newline{\footnotesize Changer le statut (règle R9)} & \texttt{id} & \texttt{ChangementStatut} & 200, 400, 401, 403, 404, 409, 422 \\ \midrule
\textsc{post} & \texttt{/\allowbreak{}candidatures/\allowbreak{}\{id\}/\allowbreak{}evaluation}\newline{\footnotesize Démarrer l'évaluation (l'entreprise ouvre le dossier)} & \texttt{id} & – & 200, 401, 403, 404, 409 \\ \midrule
\textsc{post} & \texttt{/\allowbreak{}candidatures/\allowbreak{}\{id\}/\allowbreak{}decision}\newline{\footnotesize Décision de l'entreprise (étapes 7 à 9)} & \texttt{id} & \texttt{Decision} & 200, 400, 401, 403, 404, 409, 422, 503 \\ \midrule
\textsc{post} & \texttt{/\allowbreak{}candidatures/\allowbreak{}\{id\}/\allowbreak{}suivi/\allowbreak{}reprendre}\newline{\footnotesize Rejouer les étapes de suivi d'une candidature acceptée (stage, statut de l'étudiant)} & \texttt{id} & – & 200, 401, 403, 404, 409, 503 \\ \midrule
\textsc{get} & \texttt{/\allowbreak{}etudiants/\allowbreak{}\{id\}/\allowbreak{}candidatures}\newline{\footnotesize Candidatures d'un étudiant} & \texttt{id} & – & 200, 401, 422 \\ \midrule
\textsc{get} & \texttt{/\allowbreak{}etudiants/\allowbreak{}\{id\}/\allowbreak{}candidatures/\allowbreak{}actives/\allowbreak{}compte}\newline{\footnotesize Règle R6 - nombre de candidatures actives de l'étudiant} & \texttt{id} & – & 200, 401, 422 \\ \midrule
\textsc{get} & \texttt{/\allowbreak{}offres/\allowbreak{}\{id\}/\allowbreak{}candidatures}\newline{\footnotesize Candidatures reçues pour une offre} & \texttt{id} & – & 200, 401, 422 \\ \midrule
\textsc{get} & \texttt{/\allowbreak{}entreprises/\allowbreak{}\{id\}/\allowbreak{}candidatures}\newline{\footnotesize Candidatures reçues par une entreprise} & \texttt{id}, \texttt{statut?} & – & 200, 400, 401, 422 \\
\end{longtable}
}

\subsection{S5 Notifications}
Chaque notification est enregistrée puis réellement envoyée par SMTP (Mailpit). Le statut \statut{ENVOYEE} ou \statut{ECHEC} est conservé avec l'erreur éventuelle ; la réponse est 201 dans les deux cas, car la notification a bien été enregistrée. Huit types sont gérés : confirmation de dépôt, candidature rejetée, nouvelle candidature (pour l'entreprise), acceptation, refus, demande de complément, rappel de date limite, alerte au responsable.
{\small% Généré par rapport/outils/generer_tables.py à partir de rest-notifications/src/main/resources/openapi.yaml. Ne pas modifier.
\begin{longtable}{@{}>{\raggedright\arraybackslash}p{0.06\linewidth}>{\raggedright\arraybackslash}p{0.40\linewidth}>{\raggedright\arraybackslash}p{0.16\linewidth}>{\raggedright\arraybackslash}p{0.14\linewidth}>{\raggedright\arraybackslash}p{0.12\linewidth}@{}}
\caption{Points d'accès REST de S5 Notifications (base \texttt{/api/v1}, 4 opérations)}\label{tab:rest-notifications}\\
\toprule \textbf{Verbe} & \textbf{URI et rôle} & \textbf{Paramètres} & \textbf{Corps} & \textbf{Codes} \\ \midrule
\endfirsthead
\toprule \textbf{Verbe} & \textbf{URI et rôle} & \textbf{Paramètres} & \textbf{Corps} & \textbf{Codes} \\ \midrule
\endhead
\bottomrule
\endlastfoot
\textsc{post} & \texttt{/\allowbreak{}notifications}\newline{\footnotesize Envoyer une notification} & – & \texttt{NouvelleNotification} & 201, 400, 401, 403, 422, 503 \\ \midrule
\textsc{get} & \texttt{/\allowbreak{}notifications}\newline{\footnotesize Lister les notifications} & \texttt{destinataire?}, \texttt{candidatureId?}, \texttt{type?}, \texttt{limite?} & – & 200, 400, 401, 503 \\ \midrule
\textsc{get} & \texttt{/\allowbreak{}notifications/\allowbreak{}\{id\}}\newline{\footnotesize Consulter une notification, avec le corps HTML envoyé} & \texttt{id} & – & 200, 401, 404, 422 \\ \midrule
\textsc{post} & \texttt{/\allowbreak{}rappels/\allowbreak{}executer}\newline{\footnotesize Exécuter maintenant les rappels de date limite} & \texttt{jours?} & – & 200, 400, 401, 403, 503 \\
\end{longtable}
}

\subsection{S6 Suivi des stages}
La création d'un stage est idempotente : un second appel pour la même candidature renvoie le stage existant au lieu d'en créer un doublon, ce qui permet de rejouer sans risque une étape de suivi interrompue.
{\small% Généré par rapport/outils/generer_tables.py à partir de rest-suivi/src/main/resources/openapi.yaml. Ne pas modifier.
\begin{longtable}{@{}>{\raggedright\arraybackslash}p{0.06\linewidth}>{\raggedright\arraybackslash}p{0.40\linewidth}>{\raggedright\arraybackslash}p{0.16\linewidth}>{\raggedright\arraybackslash}p{0.14\linewidth}>{\raggedright\arraybackslash}p{0.12\linewidth}@{}}
\caption{Points d'accès REST de S6 Suivi des stages (base \texttt{/api/v1}, 11 opérations)}\label{tab:rest-suivi}\\
\toprule \textbf{Verbe} & \textbf{URI et rôle} & \textbf{Paramètres} & \textbf{Corps} & \textbf{Codes} \\ \midrule
\endfirsthead
\toprule \textbf{Verbe} & \textbf{URI et rôle} & \textbf{Paramètres} & \textbf{Corps} & \textbf{Codes} \\ \midrule
\endhead
\bottomrule
\endlastfoot
\textsc{post} & \texttt{/\allowbreak{}stages}\newline{\footnotesize Créer le dossier de stage d'une candidature acceptée (règle R10)} & – & \texttt{NouveauStage} & 200, 201, 400, 401, 409, 422, 503 \\ \midrule
\textsc{get} & \texttt{/\allowbreak{}stages}\newline{\footnotesize Lister les stages} & \texttt{statut?}, \texttt{etudiantId?}, \texttt{entrepriseId?} & – & 200, 400, 401, 503 \\ \midrule
\textsc{get} & \texttt{/\allowbreak{}stages/\allowbreak{}\{id\}}\newline{\footnotesize Consulter un stage (statut et dates)} & \texttt{id} & – & 200, 401, 404, 422 \\ \midrule
\textsc{patch} & \texttt{/\allowbreak{}stages/\allowbreak{}\{id\}/\allowbreak{}statut}\newline{\footnotesize Changer le statut d'un stage} & \texttt{id} & \texttt{ChangementStatut} & 200, 400, 401, 404, 409, 422 \\ \midrule
\textsc{post} & \texttt{/\allowbreak{}stages/\allowbreak{}\{id\}/\allowbreak{}comptes-rendus}\newline{\footnotesize Enregistrer un compte rendu périodique} & \texttt{id} & \texttt{NouveauCompteRendu} & 201, 400, 401, 404, 409, 422 \\ \midrule
\textsc{get} & \texttt{/\allowbreak{}stages/\allowbreak{}\{id\}/\allowbreak{}comptes-rendus}\newline{\footnotesize Lister les comptes rendus d'un stage} & \texttt{id} & – & 200, 401, 404 \\ \midrule
\textsc{post} & \texttt{/\allowbreak{}stages/\allowbreak{}\{id\}/\allowbreak{}evaluations}\newline{\footnotesize Enregistrer une évaluation (note sur 20 et commentaires)} & \texttt{id} & \texttt{NouvelleEvaluation} & 201, 400, 401, 404, 409, 422 \\ \midrule
\textsc{get} & \texttt{/\allowbreak{}stages/\allowbreak{}\{id\}/\allowbreak{}evaluations}\newline{\footnotesize Lister les évaluations d'un stage} & \texttt{id} & – & 200, 401, 404 \\ \midrule
\textsc{get} & \texttt{/\allowbreak{}stages/\allowbreak{}\{id\}/\allowbreak{}historique}\newline{\footnotesize Historique des opérations d'un stage} & \texttt{id} & – & 200, 401, 404 \\ \midrule
\textsc{get} & \texttt{/\allowbreak{}etudiants/\allowbreak{}\{id\}/\allowbreak{}stages}\newline{\footnotesize Historique des stages d'un étudiant} & \texttt{id} & – & 200, 401, 422 \\ \midrule
\textsc{get} & \texttt{/\allowbreak{}candidatures/\allowbreak{}\{id\}/\allowbreak{}stage}\newline{\footnotesize Stage créé pour une candidature} & \texttt{id} & – & 200, 401, 404 \\
\end{longtable}
}

\section{Exemples d'échanges}

Les échanges suivants ont été enregistrés sur la pile en fonctionnement (mode direct) ; seule l'indentation du JSON a été ajoutée.

\lstinputlisting[language=HTTPfr, caption={Dépôt d'une candidature : soumission enregistrée, avec les étapes du processus}, label={lst:depot}]{generes/exemples/rest-soumission-enregistree.txt}

\lstinputlisting[language=HTTPfr, caption={Dépôt refusé par la règle R3 (date limite dépassée)}, label={lst:rejet}]{generes/exemples/rest-soumission-rejetee.txt}

\lstinputlisting[language=HTTPfr, caption={Annulation refusée pendant l'évaluation (règle R8)}, label={lst:erreur-409}]{generes/exemples/rest-erreur-409-r8.txt}

\lstinputlisting[language=HTTPfr, caption={Décision de l'entreprise : acceptation}]{generes/exemples/rest-decision.txt}

Les autres échanges enregistrés (consultation de la candidature, stage créé dans S6, notification enregistrée dans S5, erreurs 404 et 401) sont dans \code{rapport/generes/exemples/}.

\section{Référence OpenAPI et Swagger UI}

Les contrats complets, avec les schémas de chaque corps et des exemples pour chaque opération, sont dans le dépôt :
\begin{itemize}
  \item \code{rest-candidatures/src/main/resources/openapi.yaml} ;
  \item \code{rest-notifications/src/main/resources/openapi.yaml} ;
  \item \code{rest-suivi/src/main/resources/openapi.yaml}.
\end{itemize}
Une fois les services démarrés, chacun est consultable et testable dans Swagger UI (figure~\ref{fig:swagger}) : \url{http://localhost:8084/docs/}, \url{http://localhost:8085/docs/}, \url{http://localhost:8086/docs/}. Le bouton « Authorize » permet de saisir la clé d'API.

% TODO capture : Swagger UI de S4 (http://localhost:8084/docs/), liste des opérations dépliée sur « Processus ».
\capture{swagger-s4.png}{70mm}{Swagger UI de S4, servi par le service lui-même}{fig:swagger}

\section{Implémentation}

Chaque service REST est une application Jersey sur un serveur Grizzly. Les ressources JAX-RS ne contiennent que la traduction HTTP ; les règles sont dans une classe de service, testée sans serveur. Des filtres communs (module \code{common}) contrôlent la clé d'API et les clients autorisés, propagent l'identifiant de corrélation et journalisent chaque accès avec sa durée ; un unique \emph{ExceptionMapper} produit le corps d'erreur uniforme. Les appels sortants passent par un client HTTP commun, avec délais d'attente et traduction des erreurs du service appelé : un 404 de S3 devient une erreur métier claire dans S4, et une panne devient un 503.
